% !TeX root = ../main.tex

\chapter{Implementation}

\section{Prototype Overview}

The prototype is implemented in Python and built on top of Slither. The main
entry point is a price manipulation analyzer that accepts a Solidity file or
project directory, initializes Slither, constructs source and sink managers, and
runs a project-level taint tracker. The prototype also includes command-line
drivers for running taint analysis, generating representative paths, previewing
LLM prompts, and executing the staged LLM workflow.

The implementation is organized around three major packages. The
\texttt{taint\_analysis} package contains the static-analysis engine, source and
sink definitions, path tracking, path reconstruction, source canonicalization,
prompt generation, and Slither integration. The \texttt{llm\_analysis} package
contains the staged LLM engine and prompt templates. The dataset and script
directories contain audit-finding collection, repository cloning, build
verification, and dataset validation utilities.

\section{Analyzer Orchestration}

The main analyzer object is \texttt{PriceManipulationAnalyzer}. It receives a
contract path and an optional taint-analysis configuration. It then initializes
three core components:

\begin{itemize}
  \item \texttt{SlitherAnalyzer}, which parses the target project and exposes
        contracts, functions, CFG nodes, and SlithIR operations.
  \item \texttt{SourceSinkManager}, which manages taint source and sink
        definitions.
  \item \texttt{TaintTracker}, which performs source identification,
        propagation, sink matching, and path reconstruction.
\end{itemize}

The analyzer returns a dictionary containing the target path, total vulnerable
path count, path summaries, serialized paths, number of sources found, and
number of sinks found. Each serialized path includes source file and line,
source contract and function, source type, canonical source ID, sink file and
line, sink kind, sink operand selector, tainted variable, contract, and
function.

\section{Source and Sink Configuration}

The prototype defines configurable patterns for common price sources and sinks.
Oracle patterns include function names such as \texttt{latestRoundData},
\texttt{getPrice}, \texttt{getReserves}, \texttt{consult}, and
\texttt{fetchPrice}. Flash-loan patterns include \texttt{flashLoan},
\texttt{borrow}, and \texttt{executeOperation}. Price calculation sinks include
swap and quote helpers, while collateral and lending sinks include liquidation,
health factor, borrowing, depositing, withdrawing, minting, burning, and
repayment patterns.

The configuration also supports storage-taint sources based on price-like state
variable names such as \texttt{price}, \texttt{rate}, \texttt{oracle}, and
\texttt{feed}. During initialization, the analyzer scans state variables in the
target contracts and constructs a storage-read source for matched variables.

\section{Taint Engine}

The \texttt{TaintTracker} implements the core static analysis. It maintains a
taint map, a path tracker, source records, raw source candidates, sink matches,
return summaries, call contexts, and worklist metrics. The engine operates at
the project level rather than only within a single contract. This is necessary
because DeFi protocols often split price logic, accounting logic, and token
operations across multiple contracts and libraries.

The engine identifies raw source candidates and converts them into canonical
source records. A source record stores a canonical ID, display label, caller
source label, source type, exact seed IDs, provenance chain, terminal external
sites, source node, source operation, and related caller metadata. This design
allows a source wrapper to remain connected to the underlying price read.

Propagation is delegated to operation-specific handlers. These handlers process
assignments, binary operations, calls, type conversions, member accesses, tuple
unpacking, and returns. The implementation also maintains return summaries so
that a tainted return value can be propagated from a callee back to caller
contexts. This is important for helper functions that wrap price reads.

\section{Path Tracking and Grouping}

The \texttt{PathTracker} stores vulnerable paths and source aliases. It maps
exact variable identifiers to canonical source IDs, tracks reporting seed IDs,
and records source realizations created during propagation. This separation
between exact IDs and canonical IDs is important because Solidity compilation
and Slither SSA can produce many local identifiers for one semantic source.

Raw vulnerable paths are grouped into representative paths before LLM analysis.
The grouping layer reduces duplicate paths, merges sink-function records, and
preserves marked sink lines. The resulting prompt entries contain source
records, sink groups, sink function code, state writes, external calls,
representative path counts, raw path counts, and source-sink chains.

This abstraction is a practical requirement. In preliminary runs, one project
can produce hundreds of raw paths. Passing all of them directly to an LLM would
increase cost and would also make the final report difficult to interpret.

\section{Prompt Generation}

The prompt generator renders concise evidence from representative paths. It
normalizes whitespace, truncates long statements, formats source provenance,
extracts line labels, and renders source-to-sink chains. It also marks sink
lines in function bodies so that the LLM can inspect the relevant code without
receiving unrelated files.

The generated prompt JSON is saved under a project-specific output directory.
This makes each analysis run auditable: the final LLM report can be traced back
to the exact prompt entries and static-analysis summary used for that run.

\section{LLM Engine}

The \texttt{LLMEngine} orchestrates four named stages:
\texttt{classify\_canonical\_sources},
\texttt{classify\_sources\_with\_context},
\texttt{triage\_sink\_group\_impact}, and
\texttt{judge\_root\_cause\_findings}. The engine supports preview mode, which
prints prompts without calling an API, and run mode, which calls the configured
OpenAI model and writes a structured JSON report.

Every stage validates the model response. Source classification must return one
result for every expected source label and must use one of the allowed source
classes and decisions. Sink triage must return one result for every expected
sink group ID. Root-cause judgment must return one result for every expected
canonical source ID and must refer only to known evidence IDs.

This validation does not prove that the model is correct, but it prevents many
format and traceability failures. It also makes the report suitable for
automatic matching against ground truth.

\section{Dataset Utilities}

The prototype includes scripts for collecting and organizing audit findings.
The dataset CSV records report date, audit firm, severity, protocol, title,
tags, year, compilation status, source-code link, Solodit URL, GitHub link,
summary, manipulation label, source-code availability, and remarks. This format
supports filtering from a broad collected set to a narrower build-verified
evaluation set.

A separate build-manifest verifier checks curated HIGH-severity projects. It
validates project IDs, roots, toolchains, build commands, compiler versions,
required environment fields, lockfiles, and build output directories. This is
important because a finding with a public repository is not automatically
usable for static analysis. The project must be available in a reproducible
form that Slither can analyze.

\section{Execution Modes}

The main taint-analysis test driver can run in three LLM modes:

\begin{itemize}
  \item \texttt{off}: run taint analysis and generate artifacts only.
  \item \texttt{preview}: render the staged LLM prompts without API calls.
  \item \texttt{run}: execute the configured LLM workflow and write a report.
\end{itemize}

The LLM engine also has a standalone test driver. Given an existing
\texttt{llm\_prompts} JSON file, it can preview prompts, run the workflow, or
step through prompts interactively. This separation makes it possible to debug
static analysis and LLM reasoning independently.

\section{Current Implementation Limitations}

The current implementation is a prototype. It focuses on taint analysis,
representative path generation, and staged LLM report generation. Sanity
verification rules are not yet fully implemented. In particular, checks for
access control, TWAP adequacy, staleness validation, slippage protection,
fee-on-transfer behavior, and time-based controls need to be formalized.

The current sink definitions also require expansion. Preliminary artifacts show
cases where canonical sources are found but no sink sites are retained. This
can happen when the analyzed target does not include the relevant code path, or
when the sink model does not cover the protocol's specific economic operation.
These limitations are treated as part of the evaluation rather than hidden.
